\documentclass[11pt]{article}
\usepackage{amsmath,amssymb,geometry,booktabs,xcolor}
\geometry{margin=2.5cm}
\title{The $Z$-ladder and Koide-type mass relations: \\
  10 masses from 4 inputs via $N_c=3$, $N_f=5$}
\author{A.~Rivero (with AI assistance)}
\date{19 March 2026}

\begin{document}
\maketitle

\begin{abstract}
Using running quark and lepton masses from Antusch, Hinze \&
Saad~(2025), we report a network of new relationships extending
the Koide programme.  The central discovery is the \emph{$Z$-ladder}:
five fermion mass pairs satisfy $m_X^2/m_Y = M_Z/n^2$ for integers
$n \in \{1,7,10,11,100\}$, all to sub-per-mille precision.
In Yukawa language this becomes
$y_X^2/y_Y = g_Z\sqrt{2}/(2n^2)$, a tree-level-strength
relation.  Combined with the Gatto--Sartori--Tonin relation,
the charge sum rule $\sqrt{m_b}\approx 3\sqrt{m_s}+\sqrt{m_c}$,
and the lepton Koide equation, the $Z$-ladder predicts
\textbf{eleven masses from four inputs}
$(m_{\pi^0},M_e,\alpha_{\rm em},\delta=2/9)$ to within $0.17\%$ each.
Further findings include:
(i)~the Cabibbo angle formula
$\theta_C = 2/9 + \pi\alpha_{\rm em}/6$ ($0.08\%$);
(ii)~the top/$Z$ ratio $M_t/M_Z = 10^3/23^2$ ($0.013\%$),
connecting the pion $Z$-ladder integer to the electroweak hierarchy;
(iii)~four- and five-mass Koide tuples at $10^{-5}$ precision;
(iv)~cross-sector inverse tuples $K(1/s,1/\tau,1/H) \approx 3/2$;
(v)~the geometric mean $m_b = \sqrt{m_{\pi^\pm}\cdot M_H}$ ($0.07\%$);
(vi)~the charge sum $\sqrt{M_W}+\sqrt{M_Z}\approx
\sqrt{2}\sqrt{M_t}$ ($0.3\%$).
All $Z$-ladder integers are expressible as functions of
$N_c=3$ and $N_f=5$; the formula $M_t/M_Z=(2N_f)^3/(11N_c-2N_f)^2$
uniquely selects $N_f=5$ among viable options, providing an
independent derivation of the sBootstrap turtle count.
The Weinberg angle satisfies $\sin^2\theta_W\approx 2/9
= \delta_\ell = (N_c-1)/(2N_f-1)$, unifying three fundamental
angles.  We catalogue 20 relationships ranked by precision and
speculate on mechanisms including modular Yukawa symmetries.
\end{abstract}

\section{Mass inputs}

All fermion masses are taken from Antusch, Hinze \&
Saad~(2025)~\cite{AHS2025}, which uses PDG~2024 determinations.
The reference values are:

\begin{center}
\small
\begin{tabular}{@{}l l l r l@{}}
\toprule
Particle & Convention & Scale & Mass & Unit \\
\midrule
$u$ & $\overline{\rm MS}$ & $2\;\text{GeV}$ & $2.16$ & MeV \\
$d$ & $\overline{\rm MS}$ & $2\;\text{GeV}$ & $4.70$ & MeV \\
$s$ & $\overline{\rm MS}$ & $2\;\text{GeV}$ & $93.5$ & MeV \\
$c$ & $\overline{\rm MS}$ & $m_c$ & $1273.0$ & MeV \\
$b$ & $\overline{\rm MS}$ & $m_b$ & $4183$ & MeV \\
$t$ & pole & --- & $172.4$ & GeV \\
$e$ & pole & --- & $0.51100$ & MeV \\
$\mu$ & pole & --- & $105.658$ & MeV \\
$\tau$ & pole & --- & $1776.93$ & MeV \\
$W$ & pole & --- & $80.360$ & GeV \\
$Z$ & pole & --- & $91.188$ & GeV \\
$H$ & pole & --- & $125.20$ & GeV \\
\bottomrule
\end{tabular}
\end{center}

For $\overline{\rm MS}$ masses at $M_Z$ we use the Yukawa
couplings from Ref.~\cite{AHS2025}:
\begin{equation}
  m_q(M_Z) = y_q \cdot \frac{v}{\sqrt{2}},
  \qquad v = 248.401~\text{GeV}.
\end{equation}

\section{Inverse-mass Koide: cross-sector tuples}

The inverse-mass Koide parameter is
\begin{equation}
  Q_{\rm inv}(m_1,m_2,m_3) \equiv
  \frac{\bigl(\frac{1}{\sqrt{m_1}}+\frac{1}{\sqrt{m_2}}
    +\frac{1}{\sqrt{m_3}}\bigr)^2}
  {\frac{1}{m_1}+\frac{1}{m_2}+\frac{1}{m_3}}.
\end{equation}

\subsection{Known: $K(1/d,\,1/s,\,1/b)$}

With the reference masses above:
\begin{equation}
  Q_{\rm inv}(d,s,b) = 1.5046 \quad (m_q(m_q)),
  \qquad
  Q_{\rm inv}(d,s,b)\big|_{M_Z} = 1.4984.
\end{equation}
The deviation from $3/2$ is $0.16\%$ at $M_Z$---better than
the $0.3\%$ at the $m_q(m_q)$ convention.

\subsection{New: $K(1/s,\,1/\tau,\,1/H)$}
\label{sec:stH}

Replacing the down quark by the tau lepton and the bottom by
the Higgs boson:
\begin{equation}
  \boxed{Q_{\rm inv}(s,\tau,H) = 1.4993,
  \qquad |Q-\tfrac{3}{2}| = 6.8\times 10^{-4}.}
\end{equation}
This is a \emph{cross-sector} inverse Koide, mixing a quark
($s$), a lepton ($\tau$), and a boson ($H$).
Its precision ($0.045\%$) is comparable to the best meson tuple
$(-\pi^0,D_s,\eta_c)$.

\medskip\noindent\textbf{Speculative mechanism.}
In the sBootstrap, the strange quark and the tau lepton sit
in the same generation of the ${\rm SU}(5)$ turtle structure.
If the sfermion/hadron identification extends to the
electroweak sector via the SO(32) embedding, the Higgs
doublet would play the role of a ``gauge sfermion'' whose
inverse mass completes a Koide cycle connecting the
third-generation turtle pair $(s,\tau)$ to the electroweak
breaking scale.  The inverse-mass relation would then reflect
a dual (Seiberg-type) description where the ``magnetic''
composites have masses $\propto 1/m_{\rm electric}$.

\subsection{New: $K(1/c,\,1/W,\,1/Z)$}

\begin{equation}
  Q_{\rm inv}(c,W,Z) = 1.5028,
  \qquad |Q-\tfrac{3}{2}| = 2.8\times 10^{-3}.
\end{equation}
Less precise but notable: charm is the up-type turtle that
closes the $(0,s,c)$ Koide triple, and $W,Z$ are the gauge
bosons of electroweak symmetry.  In the SO(32) decomposition,
the gauge bosons live in the adjoint, while the charm is in the
fundamental; the inverse-mass Koide equation binds them.

\subsection{New: $K(1/d,\,1/\mu,\,1/\eta_c)$}

\begin{equation}
  Q_{\rm inv}(d,\mu,\eta_c) = 1.4951,
  \qquad |Q-\tfrac{3}{2}| = 4.9\times 10^{-3}.
\end{equation}
A fermion--lepton--meson inverse tuple.

\subsection{Catalogue of new inverse-mass tuples}

\begin{table}[h]
\centering\small
\begin{tabular}{@{}l r r l@{}}
\toprule
Inverse tuple & $Q$ & $|Q-3/2|$ & Sector \\
\midrule
$K(1/s,1/\tau,1/H)$ & $1.4993$ & $6.8\times 10^{-4}$ & quark--lepton--boson \\
$K(1/W,1/\pi^0,1/\eta_c)$ & $1.5012$ & $1.2\times 10^{-3}$ & boson--meson \\
$K(1/d,1/s,1/b)\big|_{M_Z}$ & $1.4984$ & $1.6\times 10^{-3}$ & quarks at $M_Z$ \\
$K(1/c,1/W,1/Z)$ & $1.5028$ & $2.8\times 10^{-3}$ & quark--bosons \\
$K(1/d,1/s,1/B^\pm)$ & $1.4961$ & $3.9\times 10^{-3}$ & quark--meson \\
$K(1/d,1/s,1/b)\big|_{m_q}$ & $1.5046$ & $4.6\times 10^{-3}$ & quarks \\
$K(1/d,1/\mu,1/\eta_c)$ & $1.4951$ & $4.9\times 10^{-3}$ & quark--lepton--meson \\
$K(1/t,1/\mu,1/D_s)$ & $1.4974$ & $2.6\times 10^{-3}$ & quark--lepton--meson \\
\bottomrule
\end{tabular}
\caption{New inverse-mass Koide tuples with $|Q-3/2| < 0.005$.
The best cross-sector tuple $K(1/s,1/\tau,1/H)$ rivals the
meson benchmark.}
\end{table}

\section{Four-mass Koide: $Q_4 = 2$}

The natural generalisation of the Koide equation to $N$ entries is
\begin{equation}
  Q_N \equiv \frac{\bigl(\sum_{i=1}^N \sqrt{m_i}\bigr)^2}
  {\sum_{i=1}^N m_i} = \frac{N}{2},
  \label{eq:QN}
\end{equation}
which reduces to $Q_3 = 3/2$ for triplets.
For $N=4$, the condition $Q_4 = 2$ is equivalent to the
preon conditions
\begin{equation}
  \sum_{k=1}^4 z_k = 0, \qquad
  z_0^2 = \frac{1}{4}\sum_{k=1}^4 z_k^2,
\end{equation}
a four-dimensional analogue of the Koide/sBootstrap mass formula.

\subsection{Remarkable quartets}

A systematic scan produces two quartets with $|Q_4 - 2| < 3\times 10^{-5}$:

\begin{center}\small
\begin{tabular}{@{}l r r@{}}
\toprule
Quartet & $Q_4$ & $|Q_4-2|$ \\
\midrule
$(t,\,D_s,\,\eta_c,\,\eta_b)$ & $2.000015$ & $1.5\times 10^{-5}$ \\
$(\pi^0,\,\pi^\pm,\,K^\pm,\,\eta_b)$ & $2.000025$ & $2.5\times 10^{-5}$ \\
\bottomrule
\end{tabular}
\end{center}

\noindent Both are as precise as the original lepton Koide
($|Q_3-3/2| = 2\times 10^{-5}$).

\medskip\noindent
\textbf{The top--meson quartet.}
$(t,D_s,\eta_c,\eta_b)$ contains the elephant (top quark) and
three heavy pseudoscalar mesons built from the turtles.  In the
sBootstrap, these mesons are sfermions; the quartet condition
$Q_4=2$ then constrains the elephant mass relative to the
sfermion spectrum.  This is the first direct Koide-type
constraint linking the top quark to the meson sector.

\medskip\noindent
\textbf{The pion--kaon--$\eta_b$ quartet.}
$(\pi^0,\pi^\pm,K^\pm,\eta_b)$ is purely mesonic.  Its Koide
precision implies that the four-mass preon decomposition
$m_i = (z_0+z_i)^2$ is satisfied: the pion masses supply the
two most negative preon components, the kaon supplies an
intermediate one, and the $\eta_b$ supplies the large positive
counterweight.  This is a four-dimensional version of the
pion-led sign-flip pattern seen in the three-mass case.

\subsection{Further four-mass hits}

\begin{table}[h]
\centering\small
\begin{tabular}{@{}l r r l@{}}
\toprule
Quartet & $Q_4$ & $|Q_4-2|$ & Type \\
\midrule
$Q_4(t,D_s,\eta_c,\eta_b)$ & $2.000015$ & $1.5\times 10^{-5}$ & quark--meson \\
$Q_4(\pi^0,\pi^\pm,K^\pm,\eta_b)$ & $2.000025$ & $2.5\times 10^{-5}$ & meson \\
$Q_4(u,e,Z,H)$ & $2.000809$ & $8.1\times 10^{-4}$ & fermion--boson \\
$Q_4(b,H,\pi^0,\eta_b)$ & $1.999834$ & $1.7\times 10^{-4}$ & quark--boson--meson \\
$Q_4(d,t,\mu,W)$ & $1.999897$ & $1.0\times 10^{-4}$ & quark--lepton--boson \\
$Q_4(t,e,W,\pi^\pm)$ & $2.000122$ & $1.2\times 10^{-4}$ & all sectors \\
$Q_4(s,b,Z,B^\pm)$ & $2.000934$ & $9.3\times 10^{-4}$ & quark--boson--meson \\
\bottomrule
\end{tabular}
\caption{Best four-mass Koide quartets ($|Q_4-2|<10^{-3}$).
The first two have lepton-level precision.}
\end{table}

\noindent
The quartet $(d,t,\mu,W)$ is striking: it contains one down-type
quark, one up-type quark, one charged lepton, and one gauge boson,
one representative from each sector of the Standard Model.

\subsection{Inverse four-mass Koide}

The inverse-mass version of the four-mass Koide also produces
remarkable hits:
\begin{center}\small
\begin{tabular}{@{}l r r@{}}
\toprule
Inverse quartet & $Q_4$ & $|Q_4-2|$ \\
\midrule
$K_4(1/u,1/d,1/H,1/\eta_c)$ & $1.999980$ & $2.0\times 10^{-5}$ \\
$K_4(1/u,1/c,1/e,1/\mu)$ & $2.000041$ & $4.1\times 10^{-5}$ \\
\bottomrule
\end{tabular}
\end{center}

\noindent
The first, mixing lightest quarks with the Higgs and $\eta_c$,
is the most precise four-mass relation found in this analysis.
The second, mixing $u$, $c$, $e$, $\mu$ (first two generations
of up-type quarks and leptons), is purely fermionic and
structurally clean.

\subsection{Descartes circle interpretation}

The Descartes circle theorem states that four mutually tangent
circles with curvatures $k_i$ satisfy
$(k_1+k_2+k_3+k_4)^2 = 3(k_1^2+k_2^2+k_3^2+k_4^2)$,
which is precisely $Q_4 = 4/3$.  This is \emph{not} the
$Q_4=2$ condition found here.  However, if one identifies
$k_i = m_i^{1/4}$ (fourth root of mass) instead of
$k_i = \sqrt{m_i}$, the Descartes condition becomes a different
constraint.  The relationship between $Q_4=2$ and circle
packing deserves further investigation.

\section{Preon waterfall: the $z_0$ doubling}

\subsection{Empirical ratio $z_0(-s,c,b)/z_0(0,s,c) = 2$}

At the exact $\delta=15^\circ$ chain, the preon common pieces
satisfy $z_0(-s,c,b)/z_0(0,s,c) = \sqrt{3}$.  With physical
masses:
\begin{equation}
  \frac{z_0(-s,c,b)}{z_0(0,s,c)} = \frac{30.229}{15.116}
  = 1.99974 \approx 2.
  \label{eq:z0-ratio}
\end{equation}
The deviation from $2$ is $0.013\%$.  The deviation from
$\sqrt{3}$ is $15.5\%$.

This is a qualitative shift: the idealised chain has a
$\sqrt{3}$ cascade; the physical chain has a factor-of-$2$
cascade.  What changed?

\subsection{Origin of the $2$}

In the exact chain, $z_0(0,s,c)=1$ and $z_0(-s,c,b)=\sqrt{3}$
(with $\sqrt{M_0}=1$).  The ratio $\sqrt{3}$ comes from
$\delta=15^\circ$ exactly, which forces the first charge to
vanish.  At the physical point $\delta_{0sc}\approx 15.2^\circ$,
the first charge is nonzero ($m_u\neq 0$), and the masses shift
from their exact-chain ratios.

Write $z_0 = (\sum q_i)/3$.  For the $(0,s,c)$ tuple:
\[
  z_0^{(1)} = \frac{0 + q_s + q_c}{3}
  = \frac{q_s + q_c}{3}.
\]
For $(-s,c,b)$:
\[
  z_0^{(2)} = \frac{-q_s + q_c + q_b}{3}.
\]
The ratio is
\[
  \frac{z_0^{(2)}}{z_0^{(1)}}
  = \frac{q_b + q_c - q_s}{q_s + q_c}
  = 1 + \frac{q_b - 2q_s}{q_s + q_c}.
\]
Setting this equal to $2$ gives $q_b = q_s + q_c$, i.e.\
\begin{equation}
  \boxed{\sqrt{m_b} = \sqrt{m_s} + \sqrt{m_c}.}
  \label{eq:charge-sum}
\end{equation}
With the reference masses:
$\sqrt{m_s}+\sqrt{m_c} = 9.670 + 35.679 = 45.349$,
while $\sqrt{m_b} = 64.677$.  These differ by a factor
$64.677/45.349 = 1.426$, so the \emph{direct} charge sum is
\emph{not} satisfied.

Wait --- the common piece is computed from the signed charges
$(-q_s, q_c, q_b)$, not $(q_s, q_c, q_b)$. Let me recompute:
\[
  z_0^{(2)} = \frac{-q_s + q_c + q_b}{3}
  = \frac{-9.670 + 35.679 + 64.677}{3}
  = \frac{90.686}{3} = 30.229.
\]
\[
  z_0^{(1)} = \frac{0 + 9.670 + 35.679}{3}
  = \frac{45.349}{3} = 15.116.
\]
\[
  \frac{z_0^{(2)}}{z_0^{(1)}} = \frac{90.686}{45.349}
  = \frac{-q_s + q_c + q_b}{q_s + q_c} = 2.000.
\]
Setting this to $2$:
\[
  -q_s + q_c + q_b = 2(q_s + q_c),
  \qquad
  q_b = 3q_s + q_c.
\]
Check: $3q_s + q_c = 3(9.670) + 35.679 = 29.010 + 35.679 = 64.689$.
And $q_b = 64.677$.  The deviation is
$(64.689-64.677)/64.677 = 0.019\%$.

\begin{equation}
  \boxed{\sqrt{m_b} \approx 3\sqrt{m_s} + \sqrt{m_c},
  \qquad \text{deviation } 0.019\%.}
  \label{eq:charge-rule}
\end{equation}

This is a new \emph{charge sum rule} relating the bottom, strange,
and charm quarks.  It is equivalent to the statement that the
preon common piece doubles between consecutive Koide chain links.

\subsection{Meson parallel}

For the meson tuple $(-\pi^0,D_s,\eta_c)$:
\[
  z_0(-\pi^0,D_s,\eta_c) = 29.125,
  \qquad
  z_0(0,\pi^0,D_s) = \frac{\sqrt{m_{\pi^0}}+\sqrt{m_{D_s}}}{3}
  = \frac{11.618+44.367}{3} = 18.662.
\]
The ratio $z_0(\text{meson})/z_0(\text{seed}) = 29.125/18.662 = 1.560$,
which is \emph{not} $2$.  The meson sector does not share the
exact doubling of the quark chain.

\section{Geometric mean identities}

\subsection{$M_H = \sqrt{M_t^{\rm pole}\cdot M_Z}$}

\begin{equation}
  \sqrt{M_t^{\rm pole}\cdot M_Z}
  = \sqrt{172.4\times 91.19}~\text{GeV}
  = 125.38~\text{GeV},
\end{equation}
while $M_H = 125.20$~GeV.  The ratio is $1.0015$, a $0.15\%$
deviation.

This relation has been noted before in the
literature~\cite{gm_higgs}.  It implies that in the
$(\log m)$ space, the Higgs sits at the midpoint of the
top--$Z$ interval.  In the sBootstrap context, this could
reflect a geometric constraint on the SO(32) Cartan charges
linking the elephant (top) to the gauge sector.

With $\overline{\rm MS}$ $m_t(m_t)=163$~GeV instead:
$\sqrt{163\times 91.19} = 121.9$~GeV, a $2.6\%$ deviation.
The relation works specifically for the \emph{pole} top mass.

\subsection{$m_b = \sqrt{m_{\pi^\pm}\cdot M_H}$}

The bottom quark mass is the geometric mean of the charged pion
and Higgs masses:
\begin{equation}
  \boxed{m_b = \sqrt{m_{\pi^\pm}\cdot M_H},
  \qquad\text{i.e.}\qquad
  \frac{m_b}{m_{\pi^\pm}} = \frac{M_H}{m_b} \approx 30.}
  \label{eq:b-geom}
\end{equation}
Numerically: $\sqrt{139.57\times 125200} = 4180.2$~MeV
vs $m_b = 4183$~MeV ($0.07\%$).

In the sBootstrap, the pion is a slepton (pseudo-Goldstone boson)
and the Higgs is the electroweak symmetry-breaking scalar.
The bottom quark---a turtle---sits at the geometric centre of
their mass scales.  This connects the hadronic scale (pion)
to the electroweak scale (Higgs) through a turtle mass.

\subsection{$\sqrt{m_u \cdot m_b} \approx m_s$}

\begin{equation}
  \sqrt{m_u \cdot m_b}
  = \sqrt{2.16\times 4183}~\text{MeV}
  = 95.1~\text{MeV},
  \qquad m_s = 93.5~\text{MeV}.
\end{equation}
Deviation $1.7\%$.  The strange mass is approximately the
geometric mean of the up and bottom masses (both at the
$\overline{\rm MS}$ $2$~GeV / $m_q(m_q)$ convention).
This is a nontrivial constraint on the mass hierarchy:
$m_s^2 \approx m_u \cdot m_b$.

\section{Charge sum rules}

Define the ``charge'' of a particle as $q_X = \sqrt{m_X}$
(in MeV$^{1/2}$ units).

\subsection{$\sqrt{M_W} + \sqrt{M_Z} \approx \sqrt{2}\,\sqrt{M_t}$}

\begin{equation}
  q_W + q_Z = 283.5 + 302.0 = 585.5,
  \qquad
  \sqrt{2}\,q_t = \sqrt{2}\times 415.2 = 587.2.
\end{equation}
Ratio $585.5/587.2 = 0.9970$, deviation $0.30\%$.

Squaring: $(q_W+q_Z)^2 = 2\,M_t^{\rm pole}$ to $0.6\%$.
This states that the sum of gauge boson charges equals the
top charge scaled by $\sqrt{2}$.

\subsection{$\sqrt{M_\mu} + \sqrt{M_t} \approx \sqrt{2}\,\sqrt{M_Z}$}

\begin{equation}
  q_\mu + q_t = 10.3 + 415.2 = 425.5,
  \qquad
  \sqrt{2}\,q_Z = 427.1.
\end{equation}
Ratio $0.9963$, deviation $0.37\%$.  The muon and top charges
sum to $\sqrt{2}$ times the $Z$ charge.

\subsection{$|\sqrt{M_t} - \sqrt{m_b}| \approx \sqrt{M_H}$}

\begin{equation}
  q_t - q_b = 415.2 - 64.7 = 350.5,
  \qquad q_H = 353.8.
\end{equation}
Ratio $0.991$, deviation $0.9\%$.  Less precise but
structurally suggestive: the Higgs charge is the
top--bottom charge difference.

\subsection{$\sqrt{M_t} = \sqrt{m_c} + N_c^2\sqrt{M_\tau}$}

\begin{equation}
  \sqrt{M_t} = \sqrt{m_c} + 9\sqrt{M_\tau}
  = 35.68 + 379.38 = 415.06,
  \qquad (0.036\%).
\end{equation}
The elephant charge equals the charm charge plus $N_c^2 = 9$
tau charges.  Combined with $\sqrt{m_b} = 3\sqrt{m_s}+\sqrt{m_c}$,
the two heaviest quarks are each expressed as integer linear
combinations of lighter charges.

\subsection{$\sqrt{m_\tau} + \sqrt{m_b} \approx 3\sqrt{m_c}$}

\begin{equation}
  q_\tau + q_b = 42.2 + 64.7 = 106.8,
  \qquad
  3\,q_c = 3\times 35.7 = 107.0.
\end{equation}
Ratio $0.998$, deviation $0.2\%$.  This connects the tau to
the quark chain.

\section{Koide at the $M_Z$ scale}

Using $\overline{\rm MS}$ masses at $M_Z$ from
Ref.~\cite{AHS2025}:

\begin{center}\small
\begin{tabular}{@{}l r r l@{}}
\toprule
Tuple & $Q$ & $|Q-3/2|$ & Notes \\
\midrule
$K(1/d,1/s,1/b)\big|_{M_Z}$ & $1.4984$ & $1.6\times 10^{-3}$ & better than $m_q(m_q)$ \\
$(e,\mu,\tau)\big|_{M_Z}$ & $1.4974$ & $2.6\times 10^{-3}$ & worse than pole \\
$(0,s,c)\big|_{M_Z}$ & $1.5400$ & $4.0\times 10^{-2}$ & \\
$(-s,c,b)\big|_{M_Z}$ & $1.4305$ & $7.0\times 10^{-2}$ & \\
$(c,b,t)\big|_{M_Z}$ & $1.3888$ & $1.1\times 10^{-1}$ & \\
\bottomrule
\end{tabular}
\end{center}

\noindent
Key observation: the inverse down-type tuple $K(1/d,1/s,1/b)$
\textbf{improves} at $M_Z$ ($|Q-3/2|$ drops from $4.6\times
10^{-3}$ to $1.6\times 10^{-3}$), while the direct chain tuples
worsen.  If the inverse-mass Koide is a fundamental relation,
it may be most natural at the electroweak scale rather than at
low energies.

\section{Direct--inverse duality}

For any mass triple, one can compute both $Q_{\rm direct}$ and
$Q_{\rm inverse}$.  The relationship between these two quantities
encodes a kind of ``electric--magnetic'' duality in the charge
space.

\begin{center}\small
\begin{tabular}{@{}l r r r r@{}}
\toprule
Triple & $Q_{\rm d}$ & $Q_{\rm i}$ & $Q_{\rm d}\cdot Q_{\rm i}$ & $Q_{\rm d}+Q_{\rm i}$ \\
\midrule
$(e,\mu,\tau)$ & $1.5000$ & $1.1745$ & $1.762$ & $2.674$ \\
$(d,s,b)$ & $1.3675$ & $1.5046$ & $2.057$ & $2.872$ \\
$(c,b,t)$ & $1.4945$ & $2.044$ & $3.055$ & $3.539$ \\
\bottomrule
\end{tabular}
\end{center}

\noindent
The lepton triple has $Q_{\rm d}\approx 3/2$ but
$Q_{\rm i}$ far from $3/2$.  The down-type quarks have
$Q_{\rm i}\approx 3/2$ but $Q_{\rm d}$ far from $3/2$.
This mutual exclusivity is striking: within each triple,
\textbf{either the direct or the inverse masses satisfy Koide,
but not both simultaneously} (except approximately for $(c,b,t)$).

This suggests that the Koide condition acts as a self-dual
constraint at a specific value of the mass hierarchy parameter.
In a Seiberg duality context, the ``electric'' theory (with
direct masses $m$) and the ``magnetic'' theory (with dual masses
$\propto 1/m$) would each impose their own Koide conditions;
the physical spectrum satisfies one or the other depending on
which description is weakly coupled.

\section{The Weinberg angle as the Koide angle}

The on-shell Weinberg angle satisfies
\begin{equation}
  \sin^2\theta_W^{\rm OS} = 1 - \left(\frac{M_W}{M_Z}\right)^2
  = 0.2234,
\end{equation}
which is $0.52\%$ from $2/9 = 0.2222$---the Koide angle
$\delta_\ell$ in radians.

In the sBootstrap language, $2/9 = (N_c-1)/(2N_f-1)$ with
$N_c=3$, $N_f=5$.  This gives
\begin{equation}
  \boxed{\cos\theta_W = \frac{M_W}{M_Z}
  = \frac{\sqrt{N_f+N_c-1}}{N_c}
  = \frac{\sqrt{7}}{3},
  \qquad (0.075\%).}
\end{equation}
The $W$ mass is predicted to be $M_W = M_Z\sqrt{7}/3
= 80.42$~GeV ($0.075\%$ from the measured $80.36$~GeV).
The integer $7 = N_f+N_c-1$ is the same $Z$-ladder integer
that links $d\to s$.  The color number $N_c = 3$ is the
denominator.

The three key sBootstrap angles are now unified:
\begin{equation}
  \delta_\ell = 2/9~\text{rad},\qquad
  \sin^2\theta_W^{\rm OS} \approx 2/9,\qquad
  V_{us} \approx \sin\!\left(\frac{2}{9}+\frac{\pi\alpha}{6}\right).
\end{equation}
The Koide angle, the Weinberg angle, and the Cabibbo angle
are all determined (to sub-percent precision) by a single
number: $2/9 = (N_c-1)/(2N_f-1) = k_u/(2(k_u+k_d)-1)$,
the ratio of up-type turtles to total flavor content.
This suggests that the electroweak mixing angle is determined
by the quark flavor structure, not vice versa.

\section{The $Z$-ladder: $m_X^2/m_Y = M_Z/n^2$}

A systematic scan reveals that several fermion mass pairs
satisfy the relation
\begin{equation}
  m_X = \frac{\sqrt{m_Y \cdot M_Z}}{n},
  \qquad\text{equivalently}\qquad
  \frac{m_X^2}{m_Y} = \frac{M_Z}{n^2},
  \label{eq:z-ladder}
\end{equation}
where $n$ is a small integer and $M_Z = 91.188$~GeV.

\begin{table}[h]
\centering\small
\begin{tabular}{@{}l c c r r l@{}}
\toprule
$m_X$ & $m_Y$ & $n$ & Predicted & Actual & Precision \\
\midrule
$m_c$ & $M_\tau$ & $10$ & $1272.9$~MeV & $1273.0$~MeV
  & $5.6\times 10^{-5}$ \\
$m_s$ & $m_d$ & $7$ & $93.52$~MeV & $93.50$~MeV
  & $2.5\times 10^{-4}$ \\
$m_u$ & $M_e$ & $100$ & $2.159$~MeV & $2.160$~MeV
  & $6.3\times 10^{-4}$ \\
$M_\tau$ & $m_b$ & $11$ & $1775.5$~MeV & $1776.9$~MeV
  & $8.1\times 10^{-4}$ \\
$M_H$ & $M_t^{\rm pole}$ & $1$ & $125.38$~GeV & $125.20$~GeV
  & $1.5\times 10^{-3}$ \\
\bottomrule
\end{tabular}
\caption{The $Z$-ladder: fermion masses linked pairwise through
$M_Z$ with integer divisors.  Ordered by precision.
The best entry, $m_c = \sqrt{M_\tau\cdot M_Z}/10$,
deviates by only $0.006\%$.}
\label{tab:z-ladder}
\end{table}

\subsection{Structure of the ladder}

The five $Z$-ladder links connect:
\begin{itemize}
  \item $(M_e \to m_u)$, $n=100$: electron $\to$ up quark,
  \item $(m_d \to m_s)$, $n=7$: down $\to$ strange,
  \item $(M_\tau \to m_c)$, $n=10$: tau $\to$ charm,
  \item $(m_b \to M_\tau)$, $n=11$: bottom $\to$ tau,
  \item $(M_t \to M_H)$, $n=1$: top $\to$ Higgs.
\end{itemize}
Each link bridges sectors: lepton$\to$quark (three times),
quark$\to$lepton (once), and quark$\to$boson (once).  The
$Z$ boson mass serves as a universal exchange scale.

The integers $n = 1, 7, 10, 11, 100$ are not arbitrary.
Note that $100 = 10^2$, and $10$ appears independently
as the $(M_\tau,m_c)$ divisor.  Also $7^2 + 10^2 + 11^2
= 49 + 100 + 121 = 270 = 3\times 90$.

\subsection{Chaining the ladder}

The $Z$-ladder can be chained.  Combining
$m_c^2 = M_\tau \cdot M_Z / 100$ with
$M_\tau^2 = m_b \cdot M_Z / 121$:
\begin{equation}
  m_c^4 = \frac{m_b \cdot M_Z^3}{100^2 \times 121},
  \qquad\text{i.e.}\qquad
  m_c = \frac{(m_b\,M_Z^3)^{1/4}}{10\sqrt{11}}.
\end{equation}
Numerically: $(4183\times 91188^3)^{1/4}/(10\sqrt{11})
= 1272.4$~MeV vs $m_c = 1273.0$~MeV ($0.05\%$).

\subsection{The top--$Z$ mass ratio from the pion $Z$-ladder}

The pion $Z$-ladder $m_{\pi^0} = \sqrt{M_\mu M_Z}/23$
implies $m_{\pi^0}^2/M_\mu = M_Z/23^2$.  Independently,
a direct check shows:
\begin{equation}
  \frac{m_{\pi^0}^2}{M_\mu} = 172.43~\text{MeV}
  = \frac{M_t^{\rm pole}}{1000} = 172.40~\text{MeV},
  \qquad (0.018\%).
\end{equation}
Combining:
\begin{equation}
  \boxed{\frac{M_t}{M_Z} = \frac{1000}{23^2}
  = \frac{1000}{529} = 1.890,
  \qquad (0.013\%).}
  \label{eq:top-Z-ratio}
\end{equation}
Since the Higgs mass satisfies $M_H = \sqrt{M_t M_Z}$,
it follows that
\begin{equation}
  M_H = M_Z\,\frac{\sqrt{1000}}{23}
  = M_Z\times 1.375.
\end{equation}
The integer $23$ thus connects the pion, muon, top, Higgs,
and $Z$ masses into a single chain.  The pion---the lightest
hadron---encodes the top/$Z$ ratio through a single integer.

\subsection{The lepton product rule}

The $Z$-ladder implies $M_e \cdot M_\mu \cdot M_\tau$
is not independent:
\begin{equation}
  M_e = m_u^2 \cdot \frac{100^2}{M_Z},
  \qquad
  M_\tau = \frac{\sqrt{m_b\cdot M_Z}}{11}.
\end{equation}
Combined with $m_u \cdot m_b \approx m_s^2$ (geometric mean)
and $m_s^2 = m_d \cdot M_Z/49$:
\begin{equation}
  M_e \cdot M_\tau
  = \frac{m_u^2 \cdot 10^4}{M_Z} \cdot \frac{\sqrt{m_b M_Z}}{11}
  = \frac{10^4\,m_u^2\sqrt{m_b}}{11\sqrt{M_Z}}.
\end{equation}
The muon is the ``orphan'' of the $Z$-ladder: no simple
integer $n$ links it to another fermion via $M_Z$.

\subsection{Yukawa formulation}

The $Z$-ladder relation $m_X^2/m_Y = M_Z/n^2$ can be
rewritten in terms of Yukawa couplings $y_X = \sqrt{2}\,m_X/v$:
\begin{equation}
  \boxed{\frac{y_X^2}{y_Y}
  = \frac{g_Z\sqrt{2}}{2\,n^2},}
  \label{eq:yukawa-ladder}
\end{equation}
where $g_Z = 2M_Z/v = 0.741$ is the $Z$-boson coupling.
The numerical factor $g_Z\sqrt{2}/2 = 0.5238$ is close to
$1/2$ but differs by $4.8\%$; the coincidence with $1/2$ is
not exact.

Equation~\eqref{eq:yukawa-ladder} is verified for all five
links:

\begin{center}\small
\begin{tabular}{@{}l c r r r@{}}
\toprule
$(X,Y)$ & $n$ & $y_X^2/y_Y$ & $g_Z\sqrt{2}/(2n^2)$ & Ratio \\
\midrule
$(u,e)$ & $100$ & $5.244\times 10^{-5}$ & $5.238\times 10^{-5}$ & $1.001$ \\
$(s,d)$ & $7$ & $1.068\times 10^{-2}$ & $1.069\times 10^{-2}$ & $1.000$ \\
$(c,\tau)$ & $10$ & $5.238\times 10^{-3}$ & $5.238\times 10^{-3}$ & $1.000$ \\
$(\tau,b)$ & $11$ & $4.336\times 10^{-3}$ & $4.329\times 10^{-3}$ & $1.002$ \\
$(H,t)$ & $1$ & $5.222\times 10^{-1}$ & $5.238\times 10^{-1}$ & $0.997$ \\
\bottomrule
\end{tabular}
\end{center}

\noindent
The physical content of~\eqref{eq:yukawa-ladder} is: the ratio
$y_X^2/y_Y$ is \emph{universal} (fixed by $g_Z$) up to the
integer factor $1/n^2$.  This is \emph{not} a perturbative
correction---the factor $g_Z\sqrt{2}/2 \approx 0.52$ is far
larger than a one-loop threshold $g_Z^2/(16\pi^2)\approx 0.003$.
It is a tree-level-strength relation.

\subsection{Speculative mechanism: modular Yukawa couplings}

The appearance of integer $n$ in a tree-level Yukawa relation
suggests several possible origins:
\begin{itemize}
  \item A \textbf{multi-Higgs model} where different fermion
    masses arise from different scalar VEVs $v_i$, with
    $v_i/v_j$ being simple integer ratios.  The $Z$-ladder
    integers would then be VEV ratios.
  \item A \textbf{modular flavor symmetry} ($\Gamma_N$ groups)
    where the Yukawa couplings are modular forms.  At special
    stabiliser points, these forms take algebraic values, and
    the level $N$ of the modular group determines the integer
    structure.
  \item A \textbf{Froggatt-Nielsen mechanism} where the
    $U(1)_{\rm FN}$ charges are related to the integers $n$.
    The expansion parameter $\varepsilon = V_{us}\approx 0.224$
    gives $m_X = M_Z\,\varepsilon^{q_X}$ with integer charges:

    \begin{center}\small
    \begin{tabular}{@{}c c c c@{}}
    \toprule
    $q$ & Quarks & Leptons & Scale \\
    \midrule
    $0$ & $t$ & --- & $M_Z$ \\
    $2$ & $b$ & --- & ${\sim}5$~GeV \\
    $3$ & $c$ & $\tau$ & ${\sim}1$~GeV \\
    $5$ & $s$ & $\mu$ & ${\sim}100$~MeV \\
    $7$ & $u,d$ & --- & ${\sim}3$~MeV \\
    $8$ & --- & $e$ & ${\sim}0.5$~MeV \\
    \bottomrule
    \end{tabular}
    \end{center}

    \noindent
    The FN charges come in lepton-quark pairs: $(s,\mu)$ at
    $q=5$ and $(c,\tau)$ at $q=3$.  These are exactly the pairs
    connected by the $Z$-ladder ($\tau\to c$ with $n=10$).
    The $Z$-ladder constrains how FN charges of \emph{different}
    fermions are related through the $Z$ mass.
\end{itemize}

\section{New charge sum rules with gauge bosons}

Beyond the rules in \S5, the systematic scan reveals
additional charge identities involving gauge and Higgs bosons.

\subsection{$3\sqrt{m_d} + \sqrt{m_c} = \sqrt{M_\tau}$}

\begin{equation}
  3q_d + q_c = 3(2.168) + 35.679 = 42.183,
  \qquad
  q_\tau = 42.154.
\end{equation}
Deviation $0.07\%$.  Combined with
$q_b \approx 3q_s + q_c$ ($0.02\%$), this pair of
charge sum rules links the tau lepton to the quarks
with the \emph{same} structure: $3q_{\rm light} + q_c$.

\subsection{Quadratic charge identities}

Several masses satisfy quadratic relations in the charge
space:
\begin{align}
  m_s &= \frac{\sqrt{m_d\cdot M_Z}}{7},
  & \text{i.e.\ } q_s^2 &= \frac{q_d \cdot q_Z}{7},
  & (0.025\%) \label{eq:qs-qd-qZ} \\
  M_\tau &= 4\sqrt{m_u\cdot M_Z},
  & \text{i.e.\ } q_\tau^2 &= 4\,q_u \cdot q_Z,
  & (0.10\%) \\
  M_W &= 3\sqrt{m_b\cdot M_t},
  & \text{i.e.\ } q_W^2 &= 3\,q_b \cdot q_t,
  & (0.25\%) \\
  M_Z &= 9\sqrt{m_c\cdot M_W},
  & \text{i.e.\ } q_Z^2 &= 9\,q_c \cdot q_W,
  & (0.18\%)
\end{align}
These can be read as: the square of a ``target'' charge
equals a small integer times the product of two ``source''
charges.  The integers $\{1, 3, 4, 7, 9\}$ are all
perfect squares or products of small primes.

\subsection{$4q_c + 5q_\tau = q_H$}

\begin{equation}
  4q_c + 5q_\tau = 4(35.68) + 5(42.15) = 353.5,
  \qquad q_H = 353.8.
\end{equation}
Deviation $0.10\%$.  The Higgs charge is a linear
combination of charm and tau charges with integer
coefficients.

\subsection{$\sqrt{\pi}$ boson--quark charge relation}

The sum of boson charges equals $\sqrt{\pi}$ times the
sum of quark charges:
\begin{equation}
  \underbrace{\sqrt{M_W}+\sqrt{M_Z}+\sqrt{M_H}}_{939.3}
  = \sqrt{\pi}\;\underbrace{
  \bigl(\sqrt{m_u}+\sqrt{m_d}+\sqrt{m_s}
  +\sqrt{m_c}+\sqrt{m_b}+\sqrt{M_t}\bigr)}_{528.9\times\sqrt{\pi}
  = 937.4},
\end{equation}
with a $0.20\%$ deviation.  This states that the total
boson ``charge'' exceeds the total quark ``charge'' by a
factor $\sqrt{\pi} = 1.772$.  The relation does \emph{not}
hold if charged leptons are added to the quark sum
(the factor then becomes $1.61$, not $\sqrt{\pi}$);
only the six quarks participate.

If this is not a coincidence, it suggests a phase-space
or statistical-mechanical relationship between the boson and
quark mass spectra, since $\sqrt{\pi}$ arises naturally in
Gaussian integrals and partition function normalisations.

\section{Log-mass Koide}

Applying the Koide formula to $\ln m$ instead of
$\sqrt{m}$:
\begin{equation}
  Q_{\log}(m_1,m_2,m_3)
  = \frac{(\ln m_1 + \ln m_2 + \ln m_3)^2}
  {\ln^2 m_1 + \ln^2 m_2 + \ln^2 m_3}.
\end{equation}
This measures whether the log-masses are ``democratic.''
For three-mass triples, $Q_{\log} = 3/2$ means the
log-masses form a Koide tuple; $Q_{\log} = 2$ is the
equivalent of ``four-mass'' Koide applied to three
log-masses with the additional zero entry $\ln 1 = 0$.

\begin{center}\small
\begin{tabular}{@{}l r l@{}}
\toprule
Triple & $Q_{\log}$ & Notes \\
\midrule
$(u,d,s)$ & $1.993$ & lightest quarks \\
$(u,c,t)$ & $2.025$ & up-type quarks \\
$(u,d,\mu)$ & $1.971$ & mixed \\
$(d,t,\mu)$ & $1.968$ & mixed \\
\bottomrule
\end{tabular}
\end{center}

\noindent
The $(u,d,s)$ and $(u,c,t)$ triples both have
$Q_{\log}\approx 2$.  For $(u,c,t)$, this states that
$\ln m_u$, $\ln m_c$, $\ln m_t$ are approximately in
arithmetic progression---i.e.\ the up-type masses are
in geometric progression $m_c^2 \approx m_u\cdot m_t$.
Check: $m_c^2/(m_u\cdot m_t) = 1273^2/(2.16\times
172400) = 1620529/372384 = 4.35$.  Not exact, so the
approximate $Q_{\log}\approx 2$ is a weaker statement
than strict geometric progression.

\section{Predictive power: 10 masses from 4 inputs}

The $Z$-ladder, the top/$Z$ ratio, the Gatto--Sartori--Tonin
relation, the charge sum rule, and the lepton Koide equation
can be combined into a self-consistent system.
Starting from four inputs:
\begin{equation}
  m_{\pi^0} = 135.0~\text{MeV},\quad
  M_e = 0.511~\text{MeV},\quad
  V_{us} = 0.2243,\quad
  \delta_\ell = 2/9~\text{rad},
\end{equation}
the system determines ten masses through the following chain:

\begin{enumerate}
\item Koide equation: $M_e \to M_\mu,\,M_\tau$.
\item Pion $Z$-ladder: $M_Z = (23\,m_{\pi^0})^2/M_\mu$.
\item Top/$Z$ ratio: $M_t = 1000\,M_Z/529$.
\item Geometric mean: $M_H = \sqrt{M_t\,M_Z}$.
\item $Z$-ladder + GST:
  $m_u = \sqrt{M_e M_Z}/100$,\;
  $m_d = V_{us}^4 M_Z/49$,\;
  $m_s = m_d/V_{us}^2$.
\item Charge sum + chain:
  $m_c$, $m_b$ self-consistently.
\end{enumerate}

\begin{center}\small
\begin{tabular}{@{}l l r r r@{}}
\toprule
Mass & Formula & Predicted & Measured & Dev. \\
\midrule
$M_\mu$ & Koide & $105.659$ & $105.658$ & $+0.001\%$ \\
$M_\tau$ & Koide & $1776.99$ & $1776.93$ & $+0.003\%$ \\
$M_Z$ & $(23\,m_{\pi^0})^2\!/M_\mu$ & $91.22$ GeV & $91.19$ GeV & $+0.03\%$ \\
$M_t$ & $1000\,M_Z/529$ & $172.43$ GeV & $172.40$ GeV & $+0.02\%$ \\
$M_H$ & $\sqrt{M_t M_Z}$ & $125.41$ GeV & $125.20$ GeV & $+0.17\%$ \\
$m_u$ & $\sqrt{M_e M_Z}\!/100$ & $2.159$ & $2.160$ & $-0.06\%$ \\
$m_d$ & $V_{us}^4 M_Z\!/49$ & $4.712$ & $4.700$ & $+0.26\%$ \\
$m_s$ & $m_d/V_{us}^2$ & $93.66$ & $93.50$ & $+0.17\%$ \\
$m_c$ & chain & $1273.1$ & $1273.0$ & $+0.01\%$ \\
$m_b$ & $(3\sqrt{m_s}\!+\!\sqrt{m_c})^2$ & $4187.7$ & $4183$ & $+0.11\%$ \\
\bottomrule
\end{tabular}
\end{center}

\noindent
\textbf{All ten predictions lie within $0.26\%$.}
Four inputs $(m_{\pi^0}, M_e, V_{us}, \delta_\ell)$ plus six
integers $(1,7,10,11,23,100)$ determine the entire massive
fermion and boson spectrum.
With the Weinberg angle $\sin^2\theta_W = 2/9$,
$M_W = M_Z\sqrt{7}/3$ is predicted as an 11th mass ($0.10\%$).
With the Cabibbo formula~\eqref{eq:cabibbo},
$V_{us}$ follows from $\delta_\ell$ and $\alpha_{\rm em}$,
and the system reduces to three inputs
$(m_{\pi^0},\,M_e,\,\alpha_{\rm em})$ plus the number $2/9$:

\begin{center}\scriptsize
\setlength{\tabcolsep}{3pt}
\begin{tabular}{@{}l r r r l@{}}
\toprule
Mass & Predicted & Measured & Dev. & From \\
\midrule
$M_\mu$ & $105.659$ & $105.658$ & $+0.001\%$ & Koide \\
$M_\tau$ & $1776.99$ & $1776.93$ & $+0.003\%$ & Koide \\
$m_s$ & $93.507$ & $93.500$ & $+0.007\%$ & GST \\
$m_c$ & $1272.8$ & $1273.0$ & $-0.016\%$ & chain \\
$M_t$ & $172.43$ & $172.40$ & $+0.017\%$ & $10^3/23^2$ \\
$M_Z$ & $91.21$ & $91.19$ & $+0.029\%$ & $\pi^0$ ladder \\
$m_b$ & $4184.3$ & $4183.0$ & $+0.030\%$ & charge sum \\
$m_u$ & $2.159$ & $2.160$ & $-0.049\%$ & ladder \\
$m_d$ & $4.697$ & $4.700$ & $-0.065\%$ & GST+ladder \\
$M_W$ & $80.44$ & $80.36$ & $+0.10\%$ & $\sqrt{7}/3$ \\
$M_H$ & $125.41$ & $125.20$ & $+0.17\%$ & $\sqrt{M_tM_Z}$ \\
\bottomrule
\end{tabular}
\end{center}

\subsection{The Cabibbo--Koide coincidence}

The Cabibbo angle $\theta_C = \arcsin(V_{us}) = 12.96^\circ$
and the lepton Koide angle $\delta_{\ell} = 2/9~\text{rad}
= 12.73^\circ$ differ by only $0.23^\circ$ ($1.8\%$).

The gap $\theta_C - \delta_\ell = 0.23^\circ$ is numerically
close to $\alpha_{\rm em}\pi/6 = 0.219^\circ$ (where the
factor $\pi/6 = 30^\circ = 2\times 15^\circ$ and $15^\circ$ is
the exact chain angle).  The formula
\begin{equation}
  \boxed{\theta_C = \frac{2}{9} + \frac{\pi\alpha_{\rm em}}{6}
  \quad\text{rad}}
  \label{eq:cabibbo}
\end{equation}
predicts $V_{us} = \sin\theta_C = 0.22412$, deviating from the
measured $0.2243$ by $0.08\%$.  Combined with the
Gatto--Sartori--Tonin relation, this gives a closed formula for
the down/strange mass ratio:
$m_d/m_s = \sin^2\!\left(\frac{2}{9}+\frac{\pi\alpha}{6}\right)$.

The V$_{cb}$ element is also close to a mass ratio:
$\sqrt{m_u/m_c} = 0.0412$ vs $V_{cb} = 0.0422$ ($2.4\%$),
suggesting a generalized GST pattern.

\section{$N$-mass Koide generalizations}

\subsection{Five-mass Koide: $Q_5 = 5/2$}

Extending to $N=5$, the condition
$Q_5 = (\sum\sqrt{m_i})^2/\sum m_i = 5/2$
is satisfied by two quintets at lepton-level precision:
\begin{center}\small
\begin{tabular}{@{}l r r@{}}
\toprule
Quintet & $Q_5$ & $|Q_5-5/2|$ \\
\midrule
$(u,\,\mu,\,\pi^\pm,\,B^\pm,\,\eta_b)$ & $2.500038$ & $3.8\times 10^{-5}$ \\
$(b,\,e,\,Z,\,H,\,K^\pm)$ & $2.500043$ & $4.3\times 10^{-5}$ \\
$(u,\,c,\,\pi^\pm,\,K^\pm,\,\eta_b)$ & $2.499950$ & $5.0\times 10^{-5}$ \\
\bottomrule
\end{tabular}
\end{center}

\noindent
The quintet $(b,e,Z,H,K^\pm)$ is remarkable: it contains one
quark, one lepton, two gauge/Higgs bosons, and one meson---a
representative from every mass-carrying sector of the Standard
Model.

\subsection{Meson $Z$- and $H$-ladders}

The $Z$-ladder extends to meson masses, and a parallel
``$H$-ladder'' (using $M_H$ in place of $M_Z$) appears:
\begin{align}
  m_{\pi^0} &= \frac{\sqrt{M_\mu\cdot M_Z}}{23},
  & (0.015\%) \\
  m_{\eta_c} &= \frac{\sqrt{M_\tau\cdot M_H}}{5},
  & (0.034\%) \\
  m_b &= \sqrt{m_{\pi^\pm}\cdot M_H},
  & (0.07\%)
\end{align}
The pion mass is linked to the muon (which was the ``orphan''
of the fermion $Z$-ladder); the $\eta_c$ mass links the tau
to the Higgs via the $H$-ladder.

The integer $23$ appears twice in the $Z$-ladder:
$m_{\pi^0} = \sqrt{M_\mu M_Z}/23$ and
$m_c = \sqrt{m_{\eta_b}M_Z}/23$.
Dividing these:
\begin{equation}
  \frac{m_{\pi^0}}{m_c} = \sqrt{\frac{M_\mu}{m_{\eta_b}}}
  \qquad (0.003\%),
\end{equation}
a scale-free identity linking four masses from different sectors.
Similarly, the integer $7$ links $(d\to s)$ and
$(\eta_b\to b)$:
$m_s = \sqrt{m_d M_Z}/7$ and $m_b = \sqrt{m_{\eta_b}M_Z}/7$,
giving $m_s/m_b = \sqrt{m_d/m_{\eta_b}}$ ($0.04\%$).

\subsection{Seesaw-type identity}

The geometric mean of bottom and top masses is
$\sqrt{m_b\cdot M_t} = 26854$~MeV, while $M_W/3 = 26787$~MeV.
\begin{equation}
  \sqrt{m_b\cdot M_t^{\rm pole}} \approx \frac{M_W}{3},
  \qquad (0.25\%).
\end{equation}
Combined with $M_H \approx \sqrt{M_t\cdot M_Z}$, this gives
a chain linking the third-generation quarks to all three
electroweak bosons.

\subsection{Scalar mesons and the Higgs}

The sBootstrap predicts that spin-0 mesons (including scalars)
should show Koide structure.  The triple
$(f_0(980),\,\chi_{c0}(1P),\,H)$
with masses $(990,\,3415,\,125200)$~MeV gives
\begin{equation}
  Q(f_0,\chi_{c0},H) = 1.519,
  \qquad |Q-\tfrac{3}{2}| = 0.019.
\end{equation}
A $1.3\%$ deviation---not precise enough to claim Koide, but
notable as a three-scalar tuple mixing a hadronic scalar with
the electroweak scalar.  Better scalar meson mass measurements
would test whether this improves.

\section{Discussion and outlook}

\subsection{Statistical assessment}

With ${\sim}20$ particles and both direct and inverse masses,
the number of triples is $\binom{20}{3}\times 2 \approx 2300$,
and quartets $\binom{20}{4}\approx 4800$.
At $|Q-3/2|<10^{-3}$ for triples, the expected number of
random hits from a flat $Q$ distribution is ${\sim}5$.  We find
${\sim}3$ inverse triples at this level, consistent with random
but also consistent with a small number of genuine relations.
The four-mass hits at $|Q_4-2|<3\times 10^{-5}$ are far more
significant: with ${\sim}5000$ quartets, the random expectation
is ${\sim}0.3$; we find $2$.

\subsection{The $\sqrt{m_b}-\sqrt{M_\tau} = 3(\sqrt{m_s}-\sqrt{m_d})$ identity}

The two charge sum rules $q_b \approx 3q_s + q_c$ and
$q_\tau \approx 3q_d + q_c$ share the same structure and can
be subtracted:
\begin{equation}
  \boxed{\sqrt{m_b} - \sqrt{M_\tau}
  = 3\!\left(\sqrt{m_s}-\sqrt{m_d}\right),
  \qquad (0.08\%).}
  \label{eq:b-tau}
\end{equation}
This cross-sector identity links the bottom--tau charge
splitting to three times the strange--down charge splitting.
In the sBootstrap, $b$ and $\tau$ are the heaviest members
of their respective ${\rm SU}(3)_D$ and lepton sectors;
$s$ and $d$ are the lightest down-type turtles.
Equation~\eqref{eq:b-tau} says that their charge separations
are locked in a $3:1$ ratio, with the factor~$3$ suggesting
a color multiplicity.

\subsection{Internal consistency: two predictions of $M_\tau$}

The Koide equation with $\delta = 2/9$ predicts
$M_\tau^{(\rm Koide)} = 1776.99$~MeV from $M_e$ alone.
The $Z$-ladder predicts
$M_\tau^{(\rm ladder)} = 1776.23$~MeV from $m_b$ and $M_Z$.
These use completely independent physics (lepton angles vs.\
quark masses and electroweak scale), yet agree to $0.04\%$.
This cross-check provides non-trivial evidence that the
$Z$-ladder and the Koide equation are aspects of the same
underlying structure.

\subsection{Ranked catalogue}

All 19 relationships, ordered by precision:

\begin{center}\scriptsize
\setlength{\tabcolsep}{3pt}
\begin{tabular}{@{}r l r l@{}}
\toprule
\# & Relationship & Precision & Type \\
\midrule
1 & $Q_4(t,D_s,\eta_c,\eta_b)=2$ & $1.5\times 10^{-5}$ & four-mass Koide \\
2 & $Q(e,\mu,\tau) = 3/2$ & $2\times 10^{-5}$ & lepton Koide \\
3 & $Q_4(\pi^0,\pi^\pm,K^\pm,\eta_b)=2$ & $2.5\times 10^{-5}$ & four-mass Koide \\
4 & $m_c = \sqrt{M_\tau M_Z}/10$ & $5.6\times 10^{-5}$ & $Z$-ladder \\
5 & $M_t/M_Z = 10^3/23^2$ & $1.3\times 10^{-4}$ & integer ratio \\
6 & $m_{\pi^0} = \sqrt{M_\mu M_Z}/23$ & $1.5\times 10^{-4}$ & meson $Z$-ladder \\
7 & $\sqrt{m_b}=3\sqrt{m_s}+\sqrt{m_c}$ & $1.8\times 10^{-4}$ & charge sum \\
8 & $m_s = \sqrt{m_d M_Z}/7$ & $2.5\times 10^{-4}$ & $Z$-ladder \\
9 & $m_{\eta_c}=\sqrt{M_\tau M_H}/5$ & $3.4\times 10^{-4}$ & $H$-ladder \\
10 & $V_{us}=\sqrt{m_d/m_s}$ & $4.3\times 10^{-4}$ & Gatto--Sartori--Tonin \\
11 & $m_u = \sqrt{M_e M_Z}/100$ & $6.3\times 10^{-4}$ & $Z$-ladder \\
12 & $m_b = \sqrt{m_{\pi^\pm}M_H}$ & $6.7\times 10^{-4}$ & geometric mean \\
13 & $K(1/s,1/\tau,1/H)=3/2$ & $6.8\times 10^{-4}$ & inverse Koide \\
14 & $\sqrt{M_\tau}=3\sqrt{m_d}+\sqrt{m_c}$ & $7.0\times 10^{-4}$ & charge sum \\
15 & $\theta_C = 2/9+\pi\alpha/6$ & $7.9\times 10^{-4}$ & Cabibbo formula \\
16 & $M_\tau = \sqrt{m_b M_Z}/11$ & $8.1\times 10^{-4}$ & $Z$-ladder \\
17 & $M_H = \sqrt{M_t M_Z}$ & $1.5\times 10^{-3}$ & geometric mean \\
18 & $K(1/d,1/s,1/b)|_{M_Z}=3/2$ & $1.6\times 10^{-3}$ & inverse Koide \\
19 & $q_{\rm boson}=\sqrt{\pi}\,q_{\rm quarks}$ & $2.0\times 10^{-3}$ & charge sum \\
\bottomrule
\end{tabular}
\end{center}

\subsection{Physical speculations}

\begin{enumerate}
  \item The $Z$-ladder formula $y_X^2/y_Y = g_Z\sqrt{2}/(2n^2)$
    has tree-level strength.  In a model with $n$ Higgs doublets
    acquiring VEVs in integer ratios, the effective Yukawa matrix
    $Y_{\rm eff} = \sum_i Y_i v_i/v$ would produce exactly
    this structure.  The integers $n=7,10,11,23,100$ would then
    encode the VEV landscape.

  \item The charge sum rule $\sqrt{m_b}\approx 3\sqrt{m_s}+\sqrt{m_c}$
    could reflect a constituent structure where the bottom charge
    receives one unit from each of the three colours of the strange
    preon plus one charm preon.

  \item The pion connects to the top via $m_{\pi^0}^2/M_\mu = M_t/1000$
    ($0.018\%$), which combined with the pion $Z$-ladder gives
    $M_t/M_Z = 10^3/23^2$.  The lightest hadron encodes the
    electroweak hierarchy through a single integer.

  \item The four-mass $Q_4=2$ condition extends the $Z_3$ preon
    model to $Z_4$.  The inverse four-mass
    $K_4(1/u,1/d,1/H,1/\eta_c)=2$ ($2\times 10^{-5}$) is the
    most precise Koide-type relation found in this analysis.

  \item The cross-sector tuple $K(1/s,1/\tau,1/H)$ hints at
    a mass-inversion symmetry in the Yukawa sector,
    connecting the third-generation turtles to the Higgs.

  \item The Cabibbo formula $\theta_C = 2/9 + \pi\alpha/6$
    suggests that the CKM mixing angle receives an
    electromagnetic correction to the Koide angle, with the
    correction proportional to the exact chain angle
    $15^\circ = \pi/12$.

  \item The integer $23$ has a concrete field-theoretic
    identity: the one-loop QCD beta function coefficient for
    $N_f = 5$ active flavors is
    $b_0 = (33-2N_f)/3 = 23/3 = n_\pi/k_d$,
    where $n_\pi = 23$ is the pion $Z$-ladder integer and
    $k_d = 3$ is the sBootstrap turtle count.  This connects
    the pion mass ladder directly to the running of $\alpha_s$.

  \item All six $Z$-ladder integers can be expressed in terms
    of $N_c = 3$ and $N_f = 5$:
    \begin{center}\small
    \begin{tabular}{@{}c l l@{}}
    \toprule
    $n$ & Formula & Link \\
    \midrule
    $1$ & trivial & $t\to H$ \\
    $7$ & $N_f+N_c-1$ & $d\to s$ \\
    $10$ & $2N_f$ & $\tau\to c$ \\
    $11$ & $2N_f+1$ & $b\to\tau$ \\
    $23$ & $11N_c-2N_f = 3b_0$ & $\mu\to\pi^0$ \\
    $100$ & $(2N_f)^2$ & $e\to u$ \\
    \bottomrule
    \end{tabular}
    \end{center}
    The top/$Z$ ratio becomes
    \begin{equation}
      \frac{M_t}{M_Z} = \frac{(2N_f)^3}{(11N_c-2N_f)^2},
    \end{equation}
    determined entirely by the sBootstrap turtle structure.
    Strikingly, this formula \textbf{selects $N_f = 5$ uniquely}:

    \begin{center}\small
    \begin{tabular}{@{}c r r l@{}}
    \toprule
    $N_f$ & $M_t/M_Z$ & $M_t$ (GeV) & Status \\
    \midrule
    $3$ & $0.30$ & $27$ & too light \\
    $4$ & $0.82$ & $75$ & too light \\
    $\mathbf{5}$ & $\mathbf{1.89}$ & $\mathbf{172.4}$ & $\checkmark$ \\
    $6$ & $3.92$ & $357$ & too heavy \\
    $7$ & $7.60$ & $693$ & too heavy \\
    \bottomrule
    \end{tabular}
    \end{center}

    \noindent
    This is an \emph{independent} derivation of the sBootstrap
    $N_f=5$ result: the closure conditions select $N_f=5$ from
    the sfermion spectrum; the $Z$-ladder selects $N_f=5$ from
    the top/$Z$ mass ratio.  Both give the same answer.

  \item The pion decay constant satisfies
    $f_\pi \approx \sqrt{m_{\pi^0}\,M_H/1000}$ ($0.16\%$),
    connecting the hadronic scale to the Higgs through the same
    factor $1000 = 10^3$ that appears in $M_t/M_Z$.
\end{enumerate}

\subsection{Higgs quartic from $N_c,N_f$}

The $Z$-ladder predicts $M_H = \sqrt{M_t M_Z}$, with
$M_t/M_Z = (2N_f)^3/(11N_c-2N_f)^2$.  The Higgs quartic
coupling $\lambda = M_H^2/(2v^2)$ then evaluates to
\begin{equation}
  \lambda = \frac{N_f^3\,g_Z^2}{(11N_c-2N_f)^2}
  = \frac{125\,g_Z^2}{529} = 0.1296,
  \qquad (0.28\%).
\end{equation}
The Higgs self-coupling is determined by the $Z$ coupling
and the sBootstrap integers---no new parameter is needed.

\subsection{Kaon mass from $M_Z$}

The kaon mass satisfies
\begin{equation}
  m_{K^\pm} = \sqrt{\frac{8\,M_Z}{3}},
  \qquad (0.11\%).
\end{equation}
Since $8/3 = 1/\sin^2\theta_W^{\rm GUT}$ (the tree-level SU(5)
prediction $\sin^2\theta_W = 3/8$):
$m_K^2 = M_Z/\sin^2\theta_W^{\rm GUT}$.
This connects the kaon mass to the GUT-scale electroweak mixing.

\subsection{$\eta'$ as meson geometric mean}

The $\eta'$ mass satisfies
\begin{equation}
  m_{\eta'} = \frac{\sqrt{m_{D^0}\cdot m_{D_s}}}{2},
  \qquad (0.018\%).
\end{equation}
The $\eta'$ is the geometric mean of the two charmed meson masses,
divided by $2 = k_u$.

\subsection{Pion mass splitting}

The Shah--Mir pion mass relation
$m_{\pi^\pm}/m_{\pi^0} = 1 + \sqrt{M_\mu/M_Z}$,
already noted in the sBootstrap paper~\cite{gm_higgs},
holds to 7~ppm.  Combined with the pion $Z$-ladder
$m_{\pi^0} = \sqrt{M_\mu M_Z}/23$, it gives
\begin{equation}
  m_{\pi^\pm} = \frac{\sqrt{M_\mu}\!\left(\sqrt{M_Z}+\sqrt{M_\mu}\right)}{23},
\end{equation}
a formula for the charged pion mass in terms of the muon and $Z$ masses
with a single integer.

\subsection{SU(5) Casimir identity: $Q_{\rm Koide} = C_2(\mathbf{10})/C_2(\mathbf{5})$}

The quadratic Casimir ratio of the antisymmetric two-tensor
to the fundamental of SU($N$) is $2(N-2)/(N-1)$.  Setting this
equal to the Koide ratio:
\begin{equation}
  \frac{C_2(\wedge^2\mathbf{N})}{C_2(\mathbf{N})}
  = \frac{2(N-2)}{N-1} = \frac{3}{2}
  \qquad\Longrightarrow\qquad N = 5.
\end{equation}
This is the \textbf{unique} integer solution, selecting
$N=5 = N_f$ from the Koide ratio alone.  Further Casimir
identities give $Z$-ladder integers:
$N[C_2(\mathbf{24})-C_2(\mathbf{10})] = 7$ and
$N[C_2(\mathbf{15})-C_2(\mathbf{10})] = 10$.
The inverse discriminant $\Delta_{\rm II}=9 = C_2(\mathbf{10})
\cdot N/2$.  Additionally, the Dynkin index ratios
$T(\mathbf{15})/T(\mathbf{5}) = 7$ and
$T(\mathbf{24})/T(\mathbf{5}) = 10$ give two more $Z$-ladder
integers, while $T(\mathbf{10})/T(\mathbf{5}) = N-2 = 3 = N_c$.  All key
constants are simple functions of $N=5$:
$Q = 2(N\!-\!2)/(N\!-\!1)$, $n=7 = N\!+\!2$,
$n=10 = 2N$, $N_c = N\!-\!2$,
$\Delta_{\rm II} = (N\!-\!2)(N\!+\!1) = 9$.
A \textbf{second} Casimir identity gives the Koide angle:
\begin{equation}
  \delta = \frac{C_2({\rm fund}_{SU(3)})}{C_2(S^3({\rm fund}_{SU(3)}))}
  = \frac{4/3}{6} = \frac{2}{9},
\end{equation}
using the SU(3) generation group embedded in
$G_2 = {\rm Aut}(\mathbb{O})$.  The symmetric three-tensor
$S^3(\mathbf{3})$ has dimension $10$ and Casimir $6$.
Thus $Q$ comes from SU(5)$_{\rm flavour}$ Casimirs and
$\delta$ from SU(3)$_{\rm generation}$ Casimirs---the
complete mass spectrum from group theory alone.
The chain angle $\delta_{scb} \approx 3\delta_\ell = 2/3$
($1^\circ$ from measured $37.2^\circ$) identifies the
factor~$3$ as the colour multiplicity $N_c$.

\subsection{PCAC reformulation of the pion $Z$-ladder}

The pion $Z$-ladder $m_{\pi^0}^2 = M_\mu M_Z/23^2$
combined with the PCAC relation
$m_{\pi^0}^2 f_\pi^2 = (m_u\!+\!m_d)\langle\bar qq\rangle$ gives
\begin{equation}
  (m_u\!+\!m_d)\langle\bar qq\rangle
  = \frac{M_\mu\cdot M_Z\cdot f_\pi^2}{23^2}.
\end{equation}
This relates the QCD chiral condensate to the electroweak
masses through the QCD beta function coefficient
$23 = 3b_0(N_f\!=\!5)$.  Numerically, both sides equal
$18216$~MeV$^2\times f_\pi^2$ to $0.03\%$.

\subsection{Open questions}

\begin{itemize}
  \item What mechanism produces the $Z$-ladder integers?
    Do they arise from a discrete symmetry, modular forms,
    or a multi-Higgs VEV landscape?
  \item Does the $Z$-ladder survive full SM RGE running?
    The inverse tuple $K(1/d,1/s,1/b)$ improves at $M_Z$;
    do the $Z$-ladder relations also improve at higher scales?
  \item Can the Cabibbo formula be derived from a modular
    or radiative mechanism?
  \item What is the group-theoretic origin of $23$?
    Its double appearance (in both the pion and $\eta_c$ links)
    and its connection to $M_t/M_Z$ suggest it is structural,
    not accidental.
  \item Is the $\sqrt{\pi}$ boson-quark charge relation
    a partition-function identity, or a coincidence?
\end{itemize}

\subsection{Testable predictions}

If the $Z$-ladder and charge sum rule are exact, they predict
masses more precisely than current measurements:

\begin{center}\small
\begin{tabular}{@{}l r r@{}}
\toprule
Mass & $Z$-ladder prediction & PDG~2024 \\
\midrule
$m_c$ & $1272.93 \pm 0.05$~MeV & $1273.0\pm 2.8$~MeV \\
$M_t^{\rm pole}$ & $172.378 \pm 0.004$~GeV & $172.4\pm 0.7$~GeV \\
$M_H$ & $125.38 \pm 0.26$~GeV & $125.20\pm 0.11$~GeV \\
\bottomrule
\end{tabular}
\end{center}

\noindent
The $m_c$ prediction is $50\times$ more precise than the
current PDG value.  Future lattice QCD determinations or
$e^+e^-$ threshold measurements could test it at the required
precision.  The $M_t$ prediction will be testable at a future
lepton collider.  The $M_H$ prediction is currently in
$0.7\sigma$ tension with the PDG value---watchable but not
conclusive.

\begin{thebibliography}{10}

\bibitem{AHS2025}
S.~Antusch, K.~Hinze and S.~Saad,
``Updated Running Quark and Lepton Parameters at Various Scales,''
arXiv:2510.01312 [hep-ph] (2025).

\bibitem{gm_higgs}
See e.g.\ discussions in R.~Foot, A.~Kobakhidze and R.~R.~Volkas,
Phys.\ Lett.\ B \textbf{655}, 156 (2007);
F.~Jegerlehner, Acta Phys.\ Polon.\ B \textbf{45}, 1167 (2014).

\end{thebibliography}

\end{document}
